\documentclass{article}
\usepackage{amsmath}
\usepackage{amssymb}

\title{Derivation of the Circle Fitting Algorithm in \texttt{fit\_arc}}
\author{Generated by Gemini}
\date{\today}

\begin{document}

\maketitle

\section{Introduction}

This document provides a step-by-step walkthrough of the calculations performed in the \texttt{fit\_arc} function, which determines the center and radius of a circle that passes through three given points in a 2D plane.

\section{Problem Statement}

Given three non-collinear points $P_1(x_1, y_1)$, $P_2(x_2, y_2)$, and $P_3(x_3, y_3)$, we want to find the center $(h, k)$ and radius $r$ of the circle that passes through these points.

\section{Derivation}

The general equation of a circle is:

$$(x - h)^2 + (y - k)^2 = r^2$$

Substituting the coordinates of the three points into this equation yields three equations:

\begin{align*}
(x_1 - h)^2 + (y_1 - k)^2 &= r^2 \\
(x_2 - h)^2 + (y_2 - k)^2 &= r^2 \\
(x_3 - h)^2 + (y_3 - k)^2 &= r^2
\end{align*}

Subtracting the second equation from the first and the third equation from the second, we get two linear equations in $h$ and $k$. Solving these equations, we find:

\begin{align*}
h &= -\frac{(x_1^2 + y_1^2)(y_3 - y_2) + (x_2^2 + y_2^2)(y_1 - y_3) + (x_3^2 + y_3^2)(y_2 - y_1)}{2(x_1(y_2 - y_3) - y_1(x_2 - x_3) + x_2y_3 - x_3y_2)} \\
k &= -\frac{(x_1^2 + y_1^2)(x_2 - x_3) + (x_2^2 + y_2^2)(x_3 - x_1) + (x_3^2 + y_3^2)(x_1 - x_2)}{2(x_1(y_2 - y_3) - y_1(x_2 - x_3) + x_2y_3 - x_3y_2)}
\end{align*}

The denominator in both expressions is the cross product of the vectors $\vec{P_1P_2}$ and $\vec{P_1P_3}$, and is calculated as:

$$cross\_product = x_1(y_2 - y_3) - y_1(x_2 - x_3) + x_2y_3 - x_3y_2$$

If $cross\_product = 0$, the points are collinear, and no unique circle can be fitted.

The numerators are determinants of matrices related to the circle center coordinates, and are calculated as:

\begin{align*}
center\_x\_determinant &= (x_1^2 + y_1^2)(y_3 - y_2) + (x_2^2 + y_2^2)(y_1 - y_3) + (x_3^2 + y_3^2)(y_2 - y_1) \\
center\_y\_determinant &= (x_1^2 + y_1^2)(x_2 - x_3) + (x_2^2 + y_2^2)(x_3 - x_1) + (x_3^2 + y_3^2)(x_1 - x_2)
\end{align*}

Thus, the center coordinates are:

\begin{align*}
h &= -\frac{center\_x\_determinant}{2 \cdot cross\_product} \\
k &= -\frac{center\_y\_determinant}{2 \cdot cross\_product}
\end{align*}

The radius $r$ is then calculated as the distance between the center $(h, k)$ and any of the three points, for example $P_1$:

$$r = \sqrt{(x_1 - h)^2 + (y_1 - k)^2}$$

\section{Implementation in \texttt{fit\_arc}}

The \texttt{fit\_arc} function implements these calculations directly. The function takes an array of two boundary points and an interior point, and returns a \texttt{CircularArc} struct representing the fitted arc, or an error if the points are collinear or the interior point is not on the arc.

\end{document}